\problemname{Birdwatchers}

The Birdwatchers' Society of San Serriffe has a curiously bloated and
ever-changing internal structure. The Society consists of $n$ chapters, and each
member of the Society belongs to exactly one chapter. The chapters are numbered
from $1$ to $n$ and the $i$'th chapter has $m_i$ members. Thus the Society has a
total of $M = m_1 + m_2 + \ldots + m_n$ members.

Each chapter is led by one of its members, who in this role is called the
\emph{officer} of the chapter. The officers are numbered the same as the
chapters, so that (for each $i = 1, \ldots, n$) officer $i$ is the one in charge
of chapter $i$.

Moreover, the officers are organized hierarchically through a system of
mentorship: each officer except one has a \emph{mentor}, who is the officer of
some other chapter. The only officer without a mentor is the President of the
Society. If officer $a$ is the mentor of officer $b$, we also say that officer
$b$ is a \emph{disciple} of officer $a$. No officer is, directly or indirectly,
a mentor of themselves; thus, by following the sequence from an officer to their
mentor, their mentor's mentor, and so on, we eventually always reach the
President.

We define the \emph{influence} of an officer as the sum of the number of members
in his chapter and of the influences of all his disciples (if he has any). It
will be easily seen that the officer with the highest influence is the
President, whose influence always equals $M$. An officer is called a
\emph{senior officer} if his influence is $\geq M/2$.

The Bylaws of the Society specify that the senior officer with the lowest
influence (amongst all the senior officers) shall act as the \emph{Treasurer} of
the Society.

From time to time, an officer (other than the President) may change his
allegiance, so that he is thenceforth the disciple of a different mentor than
before (provided that his new mentor is not one of his disciples, or his
disciples' disciples, etc.). Because of this, it can happen that the influence
of some officers changes and the role of Treasurer falls to a different officer
than before.

\section*{Task}
Write a program that reads the description of the initial state of the Society
and a sequence of changes of allegiance. Your program must print who the
Treasurer is in the initial state of the Society as well as after each change of
allegiance.

\section*{Input}
The first line contains two integers, $n$ and $q$, separated by a space; $n$ is
the number of chapters, $q$ is the number of changes of allegiance.

The next $n$ lines describe the initial state of the Society. The $i$'th of these
lines contains two integers, $s_i$ and $m_i$, separated by a space; $s_i$ is the
mentor of officer $i$ (i.e. of the officer in charge of chapter $i$), while $m_i$
is the number of members of chapter $i$. The value $s_i = 0$ indicates that
officer $i$ is the President of the Society and thus has no mentor.

The remaining $q$ lines describe the changes of allegiance. The $j$'th of these
lines contains two integers, $\hat{x}_j$ and $\hat{z}_j$, separated by a space.
The meaning of these integers is as follows. Let us denote by $t_j$ (for
$j = 0, \ldots, q$) the Treasurer after the first $j$ changes of allegiance (thus
$t_0$ is the initial Treasurer before the first change of allegiance). Then the
$j$'th change of allegiance consists of officer $z_j$ becoming the new mentor of
officer $x_j$, where $x_j = 1 + ((t_{j-1} + \hat{x}_j) \bmod n)$ and
$z_j = 1 + ((t_{j-1} + \hat{z}_j) \bmod n)$. The purpose of this representation
of the values $x_j$ and $z_j$ is to force your program to process the changes of
allegiance in the order in which they appear.

The changes of allegiance in the input data will always be valid, i.e. $z_j$ will
not be equal to $x_j$, nor will $z_j$ be a disciple of $x_j$, a disciple of a
disciple of $x_j$, etc. It is, however, possible that $z_j$ was already the
mentor of $x_j$ just before the $j$'th change (so that nothing actually changes
at that point).

Note that if your program calculates the wrong result $t_j$ at some point, it
will decode the subsequent inputs $\hat{x}_{j+1}$, $\hat{z}_{j+1}$, etc.,
incorrectly as well, and might crash with an RTE (runtime error) verdict instead
of a WA (wrong answer) verdict because the incorrectly decoded inputs might be
invalid (e.g. it might erroneously obtain a $z_{j+1}$ which is a disciple of
$x_{j+1}$).

The following limits hold:
\begin{itemize}
  \item $1 \leq n \leq 1\,000\,000$
  \item $1 \leq q \leq 30\,000$
  \item $1 \leq m_i$ for each $i = 1, \ldots, n$
  \item $m_1 + m_2 + \ldots + m_n \leq 10^9$
  \item $1 \leq \hat{x}_j \leq n$ and $1 \leq \hat{z}_j \leq n$ for each
        $j = 1, \ldots, q$.
\end{itemize}

\section*{Output}
Print the numbers $t_0, t_1, \ldots, t_q$, each on its own line, where $t_j$ is
the Treasurer after the first $j$ changes of allegiance. Naturally, each $t_j$
must be an integer from the range $1 \leq t_j \leq n$.

\section*{Scoring}
Your solution will be tested on a set of test groups.
To earn points for a group, you must pass all test cases in that group.

\noindent
\begin{tabular}{| l | l | p{12cm} |}
  \hline
  \textbf{Group} & \textbf{Points} & \textbf{Constraints} \\ \hline
  $1$ & $15$ & $n \leq 100$ \\ \hline
  $2$ & $10$ & $n \leq 1\,000$ \\ \hline
  $3$ & $50$ & $n \leq 300\,000$ \\ \hline
  $4$ & $25$ & No additional constraints. \\ \hline
\end{tabular}

\section*{Explanation of Samples}
Initially, officer $2$ is the Treasurer (hence $t_0 = 2$). In the first change of
allegiance, we read $\hat{x}_1 = 3$ and $\hat{z}_1 = 7$ and calculate
$x_1 = 1 + ((2+3) \bmod 7) = 6$ and $z_1 = 1 + ((2+7) \bmod 7) = 3$; thus,
officer $3$ becomes the new mentor of officer $6$; officer $2$ remains Treasurer
(hence $t_1 = 2$). In the second change of allegiance, we read $\hat{x}_2 = 2$
and $\hat{z}_2 = 7$ and calculate $x_2 = 1 + ((2+2) \bmod 7) = 5$ and
$z_2 = 1 + ((2+7) \bmod 7) = 3$; thus, officer $3$ becomes the new mentor of
officer $5$ and also becomes the new Treasurer (hence $t_2 = 3$).
